%! Inference and Experiments — Unit 3, Lesson 3.5 — AP Practice (Answer Key)
% MIRROR of ../ap_practice/main.tex: -key for -boxes, the header, and the correct MC option wrapped
% in \ans{}. Every \gradedblank line is identical; key mode prints each answer and display.
% NO teachernote here — guidance lives in the lesson plan.
% GRADUAL RELEASE: AP-format practice, extra credit, scored 1 to 5;
% the graded practice is the `homework` component that follows it.
% Two multiple-choice items (MC1: the scope of inference for a randomized experiment on
% volunteers — distractor (A) generalizes to all runners; MC2: reading a simulation — 46 of 200
% reshuffles reach the observed difference, so it is NOT statistically significant; distractor
% (C) is "proves no effect"), then one free-response set on green tea and blood pressure (a random
% sample, no random assignment): (a) THE CRUX in AP clothing — cause or not, with a confounder;
% (b) to whom the result generalizes; (c) the SPIRAL — design an experiment (Lesson 3.4).
% Every rubric box is full width on its own line under its part, hung with \gradedhere over
% \writelines{8} (about 2.1in, room for the drawn 2in box). Each \gradedblank line, display=
% included, is BYTE-IDENTICAL in ap_practice_key/; key mode prints the display. The
% next-lesson preview is NOT here: it closes the packet, in ../homework/.
\documentclass[12pt]{article}
\usepackage{apstats-article}
\usepackage{apstats-key}   % pulls in -boxes; do NOT also load -boxes
\usepackage{grading}
\gradingassignment{apstats-u03-l05-ap}
\gradingcourse{advanced_statistics}

\begin{document}

\pageheader{Unit 3, Lesson 3.5}{AP Practice --- Answer Key}
{\small\itshape Write each final answer in its box. Only what is inside a box is graded.}\par
% NAMESTRIP: no \namedateperiod here — mirrors the blank.

\vspace{0.06in}

% Authored BYTE-IDENTICALLY in ap_practice/ and ap_practice_key/.
\boxguard[8]
\begin{remindbox}
\small
This page is \textbf{extra credit, scored 1 to 5} --- your graded homework is the last section
of this packet. Work these AP-format reps for practice and bring questions to study hall.
\end{remindbox}

\vspace{0.06in}

\boxguard[10]
\begin{headlinebox}{goldbg}
\small
\textbf{Section I --- Multiple Choice.} Write the letter of the single best answer in its Letter box.
\end{headlinebox}

\vspace{0.06in}

\begin{enumerate}[leftmargin=1.9em, itemsep=0.10in]

  \item A sports-drink company recruits 80 runners who volunteer for a study and randomly
        assigns 40 to drink its new sports drink and 40 to drink water before a 5K race. The
        sports-drink group's mean time is faster, and the difference is statistically
        significant. Which conclusion is most appropriate?
        \begin{enumerate}[label=\textbf{(\Alph*)}, leftmargin=2.2em, itemsep=1pt]
          \item The sports drink causes faster 5K times for all runners.
          \item \ans{The sports drink causes faster 5K times for runners like these volunteers, but
                the result cannot be generalized to all runners.}
          \item The sports drink is associated with faster times, but cause and effect cannot
                be concluded because the runners were not randomly selected.
          \item The result applies to all runners, but cause and effect cannot be concluded.
          \item Nothing can be concluded, because the runners were volunteers.
                \hfill\gradedright{\textit{Letter:}~\gradedblank[id=ap1, type=choice, points=1, answer=B]{0.45in}{0.24in}}\rule[-0.2in]{0pt}{0.48in}
        \end{enumerate}

  \item A gardener randomly assigns 20 tomato plants to a new fertilizer or her usual one, 10
        each. The new-fertilizer plants produce 0.6 kg more fruit on average. In 200 random
        reshuffles of the 20 harvests, as if the fertilizer made no difference, 46 gave a
        difference of 0.6 kg or more. Which statement is best?
        \begin{enumerate}[label=\textbf{(\Alph*)}, leftmargin=2.2em, itemsep=1pt]
          \item The difference is statistically significant, because 0.6 kg is greater than 0.
          \item \ans{The difference is not statistically significant: chance alone produced a
                difference this large fairly often, so chance is a plausible explanation.}
          \item The simulation proves that the new fertilizer has no effect.
          \item The difference is statistically significant, because 46 is less than half of 200.
          \item No conclusion is possible, because the plants were not a random sample.
                \hfill\gradedright{\textit{Letter:}~\gradedblank[id=ap2, type=choice, points=1, answer=B]{0.45in}{0.24in}}\rule[-0.2in]{0pt}{0.48in}
        \end{enumerate}

\end{enumerate}

\vspace{0.08in}

\boxguard[27]
\begin{headlinebox}{goldbg}
\small
\textbf{Section II --- Free Response.} Every answer must be written \emph{in the context of the
problem} --- that is what earns credit on the AP exam.
\end{headlinebox}

\vspace{0.06in}

\begin{enumerate}[leftmargin=1.9em, itemsep=0.10in, resume]

  \item Researchers selected 1{,}000 adults at random from a list of all adult residents of a
        state. Each adult reported how many cups of green tea he or she drinks a day, and a
        nurse measured each adult's blood pressure. Adults who drink green tea every day had a
        statistically significantly lower mean blood pressure than adults who do not.

        \vspace{6pt}

        \boxguard[30]
        \textbf{(a)} Can the researchers conclude that drinking green tea \emph{causes} lower
        blood pressure? Explain, and name a variable that could confound the result.\\[4pt]
        \gradedhide{\gradedhere{\gradedblank[id=ap3a, type=rubric, points=2, display={No. The adults chose how much tea to drink; no one was randomly assigned. Tea drinkers may exercise more or eat healthier diets, which could lower blood pressure.}]{\linewidth+6pt}{\gradedlinesheight{2}}}\writelines{8}}

        \vspace{4pt}
        \boxguard[30]
        \textbf{(b)} To what group can the researchers generalize their result? Explain.\\[4pt]
        \gradedhide{\gradedhere{\gradedblank[id=ap3b, type=rubric, points=2, display={All adult residents of the state: the 1,000 were selected at random from that list. The generalized result is an association, not a cause.}]{\linewidth+6pt}{\gradedlinesheight{2}}}\writelines{8}}

        \vspace{4pt}
        \boxguard[30]
        \textbf{(c)} Describe how researchers could design an experiment to test whether
        drinking green tea causes lower blood pressure.\\[4pt]
        \gradedhide{\gradedhere{\gradedblank[id=ap3c, type=rubric, points=2, display={Randomly assign volunteers to drink green tea daily or a look-alike tea with no green tea for 8 weeks; compare the groups' mean change in blood pressure.}]{\linewidth+6pt}{\gradedlinesheight{2}}}\writelines{8}}

\end{enumerate}

\end{document}
